\documentclass[10pt]{article}

% ============================================================
% PAGE SETUP
% ============================================================
\usepackage[
    a4paper,
    left=0.45in,
    right=0.45in,
    top=0.35in,
    bottom=0.35in
]{geometry}

\usepackage{titlesec}
\usepackage{enumitem}
\usepackage{xcolor}
\usepackage{helvet}
\usepackage{tabularx}
\usepackage{ragged2e}
\usepackage{hyperref}

\renewcommand{\familydefault}{\sfdefault}
\pagestyle{empty}

\setlength{\parindent}{0pt}
\setlength{\parskip}{0pt}

% ============================================================
% HYPERLINK CONFIGURATION
% ============================================================
\hypersetup{
    colorlinks=true,
    urlcolor=black,
    linkcolor=black,
    pdfborder={0 0 0}
}

% ============================================================
% SECTION HEADINGS
% ============================================================
\titleformat{\section}
    {\normalsize\bfseries\scshape}
    {}{0em}{}
    [\titlerule]

\titlespacing*{\section}
    {0pt}
    {5pt}
    {4pt}

% ============================================================
% ENTRY HEADING
% ============================================================
\newcommand{\heading}[2]{%
    \noindent
    \textbf{#1}
    \hfill
    \textit{\small #2}
    \par
    \vspace{1pt}
}

% ============================================================
% PROJECT HEADING
% ============================================================
\newcommand{\project}[3]{%
    \noindent
    \textbf{#1}
    \hspace{2pt}
    \href{#2}{\small[\textbf{Code}]}
    \hfill
    {\small\emph{#3}}
    \par
    \vspace{1pt}
}

% ============================================================
% BULLET FORMATTING
% ============================================================
\setlist[itemize]{
    leftmargin=14pt,
    itemsep=0.5pt,
    topsep=1pt,
    partopsep=0pt,
    parsep=0pt,
    label=\textbullet
}

% ============================================================
% DOCUMENT
% ============================================================
\begin{document}

% ============================================================
% HEADER
% ============================================================
\begin{center}

{\Large\bfseries NIRANJAN KRISHNAKUMAR}\\[3pt]

{\small
Coimbatore, India
\enspace\textbullet\enspace
+91 88838 13030
\enspace\textbullet\enspace
\href{mailto:niranjankrishnakumar2005@gmail.com}
{niranjankrishnakumar2005@gmail.com}
\enspace\textbullet\enspace
\href{https://github.com/niranjan-crypt}{GitHub}
\enspace\textbullet\enspace
\href{https://niranjan-crypt.github.io}{Portfolio}
}

\end{center}

\vspace{-2pt}
\hrule
\vspace{4pt}

% ============================================================
% SUMMARY
% ============================================================
{\small\justifying
Electrical \& Computer Engineering undergraduate (B.Tech, expected 2027) applying data science and machine learning to real-world problems across computer vision, signal processing, and optimization, with one IEEE paper published and one accepted for presentation. Build end-to-end data pipelines in Python and SQL, benchmark models against classical baselines, use explainable AI (SHAP/LIME), and deliver results as interactive Streamlit apps.
\par}

% ============================================================
% EDUCATION
% ============================================================
\section{Education}

\heading
{B.Tech in Electrical and Computer Engineering}
{Expected 2027}

{\small
\textit{Amrita Vishwa Vidyapeetham, Coimbatore}
}

% ============================================================
% TECHNICAL SKILLS
% ============================================================
\section{Technical Skills}

\begin{tabularx}{\linewidth}{@{}l@{\hspace{8pt}}X@{}}

\textbf{\small Programming}
&
{\small Python, SQL, C, C++, Object-Oriented Programming (OOP), MATLAB, Bash}
\\[2pt]

\textbf{\small Data \& Apps}
&
{\small Streamlit, pandas, NumPy, Matplotlib, Data Pipelines, Feature Engineering}
\\[2pt]

\textbf{\small Traditional ML}
&
{\small Regression, Decision Trees, Random Forest, XGBoost, scikit-learn, Model Benchmarking}
\\[2pt]

\textbf{\small Deep Learning}
&
{\small PyTorch, TensorFlow, Keras, CNNs, LSTMs, Transfer Learning, Explainable AI (SHAP/LIME)}
\\[2pt]

\textbf{\small Computer Vision}
&
{\small OpenCV, Image Classification, Feature Engineering (HOG, LBP, GLCM), YOLO}
\\[2pt]

\textbf{\small Core CS \& Tools}
&
{\small Data Structures \& Algorithms, Git, Linux, Model Evaluation, MS Excel, MS PowerPoint}
\\[2pt]

\textbf{\small Systems \& Hardware}
&
{\small Raspberry Pi (Edge Inference), DSP (FFT, THD), Docker, Metaheuristics (PSO)}
\\[2pt]

\textbf{\small Cloud}
&
{\small AWS (Cloud Foundations)}

\end{tabularx}

% ============================================================
% TECHNICAL PROJECTS
% ============================================================
\section{Technical Projects}

% ------------------------------------------------------------
% PROJECT 1
% ------------------------------------------------------------
\project
{Real-Time Hyperspectral Face Identification}
{https://github.com/niranjan-crypt/Hyperspectral-Face-Identification-XAI}
{Python \textbullet\ TensorFlow \textbullet\ Raspberry Pi \textbullet\ XAI \textbullet\ Streamlit}

\begin{itemize}

\item Engineered an end-to-end spectral-spatial biometric identification
pipeline using hyperspectral imaging, mutual-information band selection,
and CNN-based deep feature representations.

\item Optimized lightweight neural inference for edge deployment on
Raspberry Pi; integrated SHAP/LIME to provide interpretable and auditable
decision verification on model predictions.

\end{itemize}

% ------------------------------------------------------------
% PROJECT 2
% ------------------------------------------------------------
\project
{Hybrid Computer Vision \& Deep Learning for Fault Diagnosis}
{https://github.com/niranjan-crypt/Hybrid-Computer-Vision-and-Deep-Learning-for-Automated-Power-Transmission-Fault-Detection}
{Python \textbullet\ TensorFlow \textbullet\ Transfer Learning \textbullet\ OpenCV \textbullet\ Streamlit}

\begin{itemize}

\item Built a multi-class defect classification pipeline for automated
power-line inspection, benchmarking transfer-learning CNNs (ResNet50,
VGG16, InceptionV3) and custom AlexNet with data augmentation against a
9-pair classical CV baseline using HOG, LBP, and GLCM.

\item Reached a 82\% accuracy, outperforming the 77.27\% Random Forest
baseline across aerial grid defect classes; applied morphological filtering
and evaluated with confusion matrices and classification reports.

\end{itemize}

% ------------------------------------------------------------
% PROJECT 3
% ------------------------------------------------------------
\project
{Physics-Informed ANN--Markov MPPT for Solar PV}
{https://github.com/niranjan-crypt/Physics-Informed-ANN-Based-MPPT-with-Markov-Chain-Based-Dynamic-Validation-for-Photovoltaic-Systems}
{Python \textbullet\ MATLAB \textbullet\ ANN \textbullet\ Stochastic Modeling}

\begin{itemize}

\item Designed an 18-15-12 ANN integrated with a Markov-chain irradiance
model for physics-informed MPPT under dynamic weather variation; scaled
the simulation dataset from 5,000 to 40,000 synthetic samples.

\item Achieved $R^2=0.9994$, $>99.5\%$ tracking efficiency, MAE = 0.08 V,
MAPE = 0.31\%, and $<5$ ms settling time in simulation; selected for
presentation as co-author at IEEE SPICES 2026.

\end{itemize}

% ------------------------------------------------------------
% PROJECT 4
% ------------------------------------------------------------
\project
{IntelliPQ: AI-Based Power Quality Monitoring \& Fault Detection}
{https://github.com/niranjan-crypt/IntelliPQ-AI-Based-Power-Quality-Monitoring-Fault-Detection}
{ TensorFlow \textbullet\ 1D-CNN + BiLSTM \textbullet\ Streamlit}

\begin{itemize}

\item Engineered an automated power-quality diagnosis pipeline using DSP
(FFT, THD, Crest Factor) and a 1D-CNN + Bi-LSTM model to classify IEEE 1159/519
grid disturbances.

\item Built a neuro-symbolic arbitration engine overriding AI conflicts with
physical safety rules; deployed real-time diagnostics and audit certificates
in Streamlit.

\end{itemize}

% ------------------------------------------------------------
% PROJECT 5
% ------------------------------------------------------------
\project
{Hybrid Optimization for Power-System Congestion Management}
{https://github.com/niranjan-crypt/Hybrid-Optimization-for-Congestion-Management-in-Deregulated-Power-Systems}
{MATLAB \textbullet\ Constrained Optimization \textbullet\ DC-OPF}

\begin{itemize}

\item Formulated a constrained DC-OPF framework on the IEEE 14-bus
transmission system to minimize active-power load shedding under line
congestion constraints.

\item Developed an adaptive hybrid metaheuristic combining Orca, Krill Herd,
and Spotted Hyena algorithms with stagnant-iteration diversification,
reducing load shedding to 3.39 MW (IEEE EPREC 2026).

\end{itemize}

% ============================================================
% INTERNSHIP & CERTIFICATIONS
% ============================================================
\section{Internship \& Certifications}

\heading
{Industry Intern --- Larsen \& Toubro (L\&T), Chennai}
{June 2026}

\begin{itemize}

\item Built simulation models of EV powertrains and battery management
systems to analyze PMSM/BLDC motor-control behavior, DC--DC converter
operation, battery integration, and CAN-based communication across
subsystems.

\item Worked with railway electrical systems and SCADA infrastructure
alongside the core powertrain modeling work.

\end{itemize}

% ------------------------------------------------------------
% AWS CERTIFICATION
% ------------------------------------------------------------
\noindent
\textbf{AWS Academy Graduate --- Cloud Foundations}
\hspace{2pt}
\href{https://www.credly.com/badges/eaed5b78-e6d1-426e-8be1-6a1c66f2bc7e/print}
{\small[\textbf{Certificate}]}
\hfill
\textit{\small Feb 2026}
\par
\vspace{1pt}

\begin{itemize}

\item Completed an AWS Academy program covering core cloud
concepts, services, security, architecture, pricing, and support;
credential issued as a verifiable Credly digital badge.

\end{itemize}

% ============================================================
% PUBLICATIONS
% ============================================================
\section{Publications}

\begin{itemize}

\item Co-author of a paper selected for IEEE SPICES 2026: ``Physics-Informed
ANN-Based MPPT with Markov-Chain Validation for Photovoltaic Systems.''

\item Co-author, ``Hybrid Optimization for Congestion Management in
Deregulated Power Systems,'' IEEE EPREC 2026.
\href{https://doi.org/10.1109/EPREC66546.2026.11412040}
{[\textbf{DOI: 10.1109/EPREC66546.2026.11412040}]}

\end{itemize}

\end{document}
